\documentclass[a4paper]{article}

\usepackage[spanish]{babel}
\usepackage{graphicx}
\usepackage{amsmath, amssymb}
\usepackage[margin=2cm]{geometry}
\usepackage{fancyhdr}
\usepackage{enumerate}
\usepackage[shortlabels]{enumitem}
\usepackage{parskip}
\usepackage[most]{tcolorbox}
\usepackage[hidelinks]{hyperref}
\usepackage{float}
\usepackage{booktabs}
\usepackage{mathrsfs}
\usepackage{tikz}
\usepackage{forest}



% cabecera
\pagestyle{fancy}
\fancyhead[l]{Mecánica celeste}
\fancyhead[c]{Notas para Quiz 5}
\fancyhead[r]{\today}
\fancyfoot[c]{\thepage}
\renewcommand{\headrulewidth}{0.2pt} % linea horizontal


\begin{document}

\begin{center}
    \huge\textbf{Demostraciones Importantes para el Quiz Final} \\
    \large Basado en el libro de Mecánica Celeste de Zuluaga
\end{center}

\tableofcontents
\newpage

\section{Llegar a la Mínima Acción desde Euler-Lagrange}

\subsection*{Desarrollo}

Partimos de las ecuaciones de Euler-Lagrange (Zuluaga, Sección 9.4):

\begin{equation}
\frac{d}{dt}\left(\frac{\partial L}{\partial \dot{q}_j}\right) - \frac{\partial L}{\partial q_j} = 0
\end{equation}

\textbf{Paso 1:} Integrar y multiplicar por $-\eta_j$ las E-L:

\begin{equation}
\int_{t_1}^{t_2} \left[\frac{d}{dt}\left(\frac{\partial L}{\partial \dot{q}_j}\right) - \frac{\partial L}{\partial q_j}\right](-\eta_j) dt = 0
\end{equation}

\textbf{Paso 2:} Integración por partes del primer término:

\begin{align}
\int_{t_1}^{t_2} \frac{d}{dt}\left(\frac{\partial L}{\partial \dot{q}_j}\right)(-\eta_j) dt &= 
\left.-\frac{\partial L}{\partial \dot{q}_j}\eta_j\right|_{t_1}^{t_2} + \int_{t_1}^{t_2} \frac{\partial L}{\partial \dot{q}_j}\dot{\eta}_j dt
\end{align}

\textbf{Paso 3:} Condiciones en los extremos (Zuluaga, Sección 9.6):

\begin{equation}
\eta(t_1) = \eta(t_2) = 0 \quad \text{(se anula en extremos)}
\end{equation}

Por lo tanto, el término de frontera se anula:

\begin{equation}
\left.-\frac{\partial L}{\partial \dot{q}_j}\eta_j\right|_{t_1}^{t_2} = 0
\end{equation}

\textbf{Paso 4:} Derivar las simetrías para despejar $\eta_j$ y reemplazar:

\begin{equation}
\int_{t_1}^{t_2} \left[\frac{\partial L}{\partial \dot{q}_j}\dot{\eta}_j + \frac{\partial L}{\partial q_j}\eta_j\right] dt = 0
\end{equation}

\textbf{Paso 5:} Identificar derivada total de $L$:

\begin{equation}
\frac{dL}{dt} = \sum_j \left(\frac{\partial L}{\partial q_j}\dot{q}_j + \frac{\partial L}{\partial \dot{q}_j}\ddot{q}_j\right) + \frac{\partial L}{\partial t}
\end{equation}

\textbf{Paso 6:} Definir la acción (Zuluaga, Ec. 9.40):

\begin{equation}
S = \int_{t_1}^{t_2} L(q_j, \dot{q}_j, t) dt
\end{equation}

La variación de la acción es:

\begin{equation}
\delta S = \int_{t_1}^{t_2} \sum_j \left(\frac{\partial L}{\partial q_j} - \frac{d}{dt}\frac{\partial L}{\partial \dot{q}_j}\right)\eta_j dt = 0
\end{equation}

Esto demuestra el \textbf{Principio de Hamilton}: la trayectoria real hace estacionaria la acción.

\section{Cuadratura de Simetrías Espaciales de Noether}

\subsection*{Desarrollo}

Según el teorema de Noether (Zuluaga, Sección 9.7), para cada simetría continua existe una cantidad conservada.

\textbf{Paso 1:} Identificar el generador de la transformación en la simetría:

Para una transformación infinitesimal:
\begin{equation}
q_j \rightarrow q_j + \varepsilon K_j(q)
\end{equation}
donde $K_j$ es el generador de la transformación.

\textbf{Paso 2:} Derivada total respecto a $\varepsilon$ de $L$:

\begin{equation}
\frac{dL}{d\varepsilon} = \sum_j \left(\frac{\partial L}{\partial q_j}\frac{\partial q_j}{\partial \varepsilon} + \frac{\partial L}{\partial \dot{q}_j}\frac{\partial \dot{q}_j}{\partial \varepsilon}\right)
\end{equation}

\textbf{Paso 3:} Derivar simetrías para despejar $K_j$ y reemplazar:

Para una simetría, $L$ es invariante bajo la transformación:
\begin{equation}
\frac{dL}{d\varepsilon} = 0
\end{equation}

\textbf{Paso 4:} Identificar el término de E-L y reemplazarlo:

Usando las ecuaciones de Euler-Lagrange:
\begin{equation}
\frac{\partial L}{\partial q_j} = \frac{d}{dt}\left(\frac{\partial L}{\partial \dot{q}_j}\right)
\end{equation}

\textbf{Paso 5:} Tenemos una derivada de producto (la cuadratura):

\begin{align}
0 &= \sum_j \left[\frac{d}{dt}\left(\frac{\partial L}{\partial \dot{q}_j}\right)K_j + \frac{\partial L}{\partial \dot{q}_j}\dot{K}_j\right] \\
&= \frac{d}{dt}\left[\sum_j \frac{\partial L}{\partial \dot{q}_j}K_j\right]
\end{align}

Por lo tanto, la \textbf{cantidad conservada} es (Zuluaga, Ec. 9.45):

\begin{equation}
P = \sum_j \frac{\partial L}{\partial \dot{q}_j}K_j = \text{constante}
\end{equation}

\section{Lagrangiano 2C en Coordenadas de Jacobi}

\subsection*{Desarrollo}

Para el problema de dos cuerpos, introducimos las coordenadas de Jacobi (Zuluaga, Sección 9.8.4):

\begin{align}
\vec{r}_{CM} &= \frac{m_1\vec{r}_1 + m_2\vec{r}_2}{m_1 + m_2} \\
\vec{r} &= \vec{r}_1 - \vec{r}_2
\end{align}

El Lagrangiano del sistema es:

\begin{equation}
L_{2D} = \frac{1}{2}M\dot{\vec{r}}_{CM}^2 + \frac{1}{2}m_r\dot{\vec{r}}^2 + \frac{GMm_r}{r}
\end{equation}

donde:
\begin{itemize}
    \item $M = m_1 + m_2$ (masa total)
    \item $m_r = \frac{m_1m_2}{m_1+m_2}$ (masa reducida)
    \item $r = |\vec{r}|$ (distancia relativa)
\end{itemize}

\textbf{Nota:} Este Lagrangiano es covariante bajo cualquier rotación.

\subsection*{¿No se usa la demostración anterior?}

En realidad, sí se puede usar. La derivación completa parte de:

\begin{equation}
L = \frac{1}{2}m_1\dot{\vec{r}}_1^2 + \frac{1}{2}m_2\dot{\vec{r}}_2^2 + \frac{Gm_1m_2}{|\vec{r}_1 - \vec{r}_2|}
\end{equation}

Sustituyendo las coordenadas de Jacobi y simplificando, se obtiene el Lagrangiano separado en movimiento del centro de masa y movimiento relativo.

\section{Invariante Temporal y Función de Jacobi}

\subsection*{Desarrollo}

\textbf{Paso 1:} Traslación temporal $t' = t + \varepsilon$, $\frac{\partial L}{\partial t} = 0$

Cuando el Lagrangiano no depende explícitamente del tiempo, existe una cantidad conservada.

\textbf{Paso 2:} Derivada total (Zuluaga, Sección 9.7.4):

\begin{equation}
\frac{dL}{dt} = \sum_j \left(\frac{\partial L}{\partial q_j}\dot{q}_j + \frac{\partial L}{\partial \dot{q}_j}\ddot{q}_j\right) + \frac{\partial L}{\partial t}
\end{equation}

\textbf{Paso 3:} Identificar término de E-L y sustituirlo:

Usando $\frac{\partial L}{\partial q_j} = \frac{d}{dt}\left(\frac{\partial L}{\partial \dot{q}_j}\right)$:

\begin{align}
\frac{dL}{dt} &= \sum_j \left[\frac{d}{dt}\left(\frac{\partial L}{\partial \dot{q}_j}\right)\dot{q}_j + \frac{\partial L}{\partial \dot{q}_j}\ddot{q}_j\right] + \frac{\partial L}{\partial t} \\
&= \frac{d}{dt}\left[\sum_j \frac{\partial L}{\partial \dot{q}_j}\dot{q}_j\right] + \frac{\partial L}{\partial t}
\end{align}

\textbf{Paso 4:} $\frac{\partial L}{\partial t}$ menos derivada de un producto, definir $h$ (Zuluaga, Ec. 9.47):

\begin{equation}
h = \sum_j \frac{\partial L}{\partial \dot{q}_j}\dot{q}_j - L
\end{equation}

Entonces:
\begin{equation}
\frac{dh}{dt} = -\frac{\partial L}{\partial t}
\end{equation}

Si $\frac{\partial L}{\partial t} = 0$, entonces $h = \text{cte}$ (función de Jacobi conservada).

\section{Invariancia Rotacional (2C)}

\subsection*{Desarrollo}

Para rotaciones en el plano, el momento angular se conserva.

\textbf{Paso 1:} Identificar generadores $K$ de la simetría:

Para una rotación infinitesimal en el plano $xy$:
\begin{align}
x &\rightarrow x + \varepsilon(-y) \\
y &\rightarrow y + \varepsilon(x)
\end{align}

Los generadores son:
\begin{equation}
K_x = -y, \quad K_y = x
\end{equation}

\textbf{Paso 2:} Aplicar cuadratura de Noether como invariantes:

\begin{align}
P &= \sum_j \frac{\partial L}{\partial \dot{q}_j}K_j \\
&= \frac{\partial L}{\partial \dot{x}}(-y) + \frac{\partial L}{\partial \dot{y}}(x) \\
&= p_x(-y) + p_y(x) \\
&= xp_y - yp_x
\end{align}

Por lo tanto, el \textbf{momento angular} se conserva:

\begin{equation}
L_z = xp_y - yp_x = \text{constante}
\end{equation}

Para el problema de dos cuerpos en coordenadas polares:

\begin{equation}
L_z = m_r r^2 \dot{\theta} = \text{constante}
\end{equation}

\section{Función de Jacobi Igual a Energía}

\subsection*{Desarrollo}

Demostraremos que $h = T + U = E$ bajo ciertas condiciones (Zuluaga, Proposición 9.5).

\textbf{Paso 1:} Partir de la definición (Zuluaga, Ec. 9.47):

\begin{equation}
h = \sum_j \frac{\partial L}{\partial \dot{q}_j}\dot{q}_j - L
\end{equation}

\textbf{Paso 2:} Reemplazar $L = T - U$:

\begin{equation}
h = \sum_j \frac{\partial (T-U)}{\partial \dot{q}_j}\dot{q}_j - (T-U)
\end{equation}

\textbf{Paso 3:} Para demostrar que es homogénea de grado 2:

Si $T$ es cuadrática en las velocidades:
\begin{equation}
T = \frac{1}{2}\sum_i m_i (\vec{v}_i \cdot \vec{v}_i)
\end{equation}

En coordenadas generalizadas:
\begin{equation}
\vec{v}_i = \sum_j \frac{\partial \vec{r}_i}{\partial q_j}\dot{q}_j
\end{equation}

\textbf{Paso 4:} Teorema de Euler para funciones homogéneas:

Si $f(\lambda x) = \lambda^k f(x)$, entonces:
\begin{equation}
\sum_j x_j \frac{\partial f}{\partial x_j} = k f
\end{equation}

Para $T$ homogénea de grado 2 en $\dot{q}_j$:
\begin{equation}
\sum_j \dot{q}_j \frac{\partial T}{\partial \dot{q}_j} = 2T
\end{equation}

\textbf{Paso 5:} Sustituir (Zuluaga, Sección 9.7.4):

\begin{align}
h &= \sum_j \frac{\partial T}{\partial \dot{q}_j}\dot{q}_j - (T-U) \\
&= 2T - T + U \\
&= T + U = E
\end{align}

Por lo tanto, \textbf{la función de Jacobi es igual a la energía total} cuando:
\begin{itemize}
    \item $T$ es homogénea de grado 2 en $\dot{q}_j$
    \item $U$ no depende de $\dot{q}_j$
    \item Las transformaciones no dependen explícitamente de $t$
\end{itemize}

\section{Constante de Jacobi}

\subsection*{Desarrollo}

La constante de Jacobi $C_J$ es fundamental en el problema restringido de los tres cuerpos (Zuluaga, Sección 8.7).

\textbf{Paso 1:} Expresar $h$ como $\sum_i \frac{\partial L}{\partial \dot{q}_i}\dot{q}_i - L$, reemplazar cada derivada:

En el sistema rotante del CRTBP con coordenadas $(x, y, z)$:

\begin{equation}
h = \frac{1}{2}(\dot{x}^2 + \dot{y}^2 + \dot{z}^2) + V_{mod}(x,y,z)
\end{equation}

donde $V_{mod}$ es el potencial modificado.

\textbf{Paso 2:} Llegamos a (Zuluaga, Ec. 8.23):

\begin{equation}
h = \frac{1}{2}(\dot{x}^2 + \dot{y}^2 + \dot{z}^2) + V_{mod} = \frac{1}{2}\vec{v}^2 + V_{mod}
\end{equation}

\textbf{Paso 3:} La constante de Jacobi (Zuluaga, Ec. 8.24):

\begin{equation}
C_J = 2\left(\frac{1-\alpha}{r_1} + \frac{\alpha}{r_2}\right) + (x^2 + y^2) - \vec{v}^2
\end{equation}

Comparando con $h$:
\begin{align}
C_J &= 2\left[\frac{1-\alpha}{r_1} + \frac{\alpha}{r_2} + \frac{1}{2}(x^2+y^2)\right] - \vec{v}^2 \\
&= 2V_{mod} - \vec{v}^2 \\
&= -2\left[\frac{1}{2}\vec{v}^2 - V_{mod}\right] \\
&= -2h
\end{align}

Por lo tanto:
\begin{equation}
\boxed{C_J = -2h}
\end{equation}

\textbf{Nota:} Nos tendrían que dar esto para usar.

\section{Ecuaciones Canónicas de Hamilton}

\subsection*{Desarrollo}

\textbf{Paso 1:} Tomar la definición de $dH$ por derivada total y por $H = \sum p_i\dot{q}_i - L$ derivada:

(A) Derivada total de $H$:
\begin{equation}
dH = \sum_j \left(\frac{\partial H}{\partial q_j}dq_j + \frac{\partial H}{\partial p_j}dp_j\right) + \frac{\partial H}{\partial t}dt
\end{equation}

(B) De la definición $H = \sum p_i\dot{q}_i - L$:
\begin{align}
dH &= \sum_j (p_j d\dot{q}_j + \dot{q}_j dp_j) - dL \\
&= \sum_j (p_j d\dot{q}_j + \dot{q}_j dp_j) - \sum_j \left(\frac{\partial L}{\partial q_j}dq_j + \frac{\partial L}{\partial \dot{q}_j}d\dot{q}_j\right) - \frac{\partial L}{\partial t}dt
\end{align}

\textbf{Paso 2:} Sustituir los momentos conjugados de la derivada total de $L$ en (B):

Recordando que $p_j = \frac{\partial L}{\partial \dot{q}_j}$:

\begin{equation}
dH = \sum_j (\dot{q}_j dp_j - \frac{\partial L}{\partial q_j}dq_j) - \frac{\partial L}{\partial t}dt
\end{equation}

\textbf{Paso 3:} Al desarrollar la derivada de producto de (B) se anulan términos con la derivada total de $-L$:

Los términos $p_j d\dot{q}_j$ y $-\frac{\partial L}{\partial \dot{q}_j}d\dot{q}_j$ se cancelan porque $p_j = \frac{\partial L}{\partial \dot{q}_j}$.

\textbf{Paso 4:} (B) se convierte en (Zuluaga, Sección 10.3):

\begin{equation}
dH = \sum_j (\dot{q}_j dp_j - \dot{p}_j dq_j) - \frac{\partial L}{\partial t}dt
\end{equation}

usando $\frac{\partial L}{\partial q_j} = \dot{p}_j$ (Euler-Lagrange).

\textbf{Paso 5:} Comparar término a término (A) y (B):

i) Coeficiente de $dt$:
\begin{equation}
\frac{\partial H}{\partial t} = -\frac{\partial L}{\partial t}
\end{equation}

ii) Coeficiente de $dq_j$:
\begin{equation}
\frac{\partial H}{\partial q_j} = -\dot{p}_j
\end{equation}

iii) Coeficiente de $dp_j$:
\begin{equation}
\frac{\partial H}{\partial p_j} = \dot{q}_j
\end{equation}

Estas son las \textbf{ecuaciones canónicas de Hamilton} (Zuluaga, Ecs. 10.10-10.11):

\begin{equation}
\boxed{
\begin{aligned}
\dot{q}_j &= \frac{\partial H}{\partial p_j} \\
\dot{p}_j &= -\frac{\partial H}{\partial q_j} \\
\frac{\partial H}{\partial t} &= -\frac{\partial L}{\partial t}
\end{aligned}}
\end{equation}

\section*{Referencias}

Todas las demostraciones están basadas en el libro \textbf{``Mecánica Celeste: Teoría, algoritmos y problemas''} de Jorge I. Zuluaga, Universidad de Antioquia, 2020.

\begin{itemize}
    \item Sección 9.4: Ecuaciones de Lagrange
    \item Sección 9.6: Principio de Hamilton
    \item Sección 9.7: Simetrías y leyes de conservación
    \item Sección 9.8: Mecánica celeste en el formalismo Lagrangiano
    \item Sección 10.3: Ecuaciones canónicas de Hamilton
    \item Sección 8.7: La constante de Jacobi
\end{itemize}

\end{document}
\bibliographystyle{plainurl}
\bibliography{bibliografia} 
\end{document}
